\section{まとめ}

本研究で得られた重要な結論は，ランサムウェア防御を
「どれだけ早く検知するか」だけで評価してはいけない，という点である．
潜伏期間が重尾的であれば，検知を強化しても，長く潜伏した感染が
バックアップ世代を汚染してから見つかる場合が残る．このとき損失を
左右するのは，検知後にどこまで古い清浄なバックアップへ戻れるかである．
したがってバックアップ運用は，障害後の保険ではなく，検知投資と並ぶ
ランサムウェア防御の中核である．

本稿の式は，バックアップ運用に対して次の設計指針を与える．
保持世代数を $n$，バックアップ間隔を $\Delta$ とすると，清浄な
復旧点へ戻れる期間は $n\Delta$ である．内点設計条件のもとで，
この釣合いから得られる目安は
\[
  \Dstar=\frac{L-(d_1\tau_R+d_0)}{d_1\kappa_s n}
\]
であり，等価的に，清浄バックアップを残す目安期間は
$n\Dstar=\{L-(d_1\tau_R+d_0)\}/(d_1\kappa_s)$ である．
全損被害と通常復旧時の被害の差が大きいほど，より古い復旧点を
残すべきである．保持世代数を増やせるなら，同じ復旧可能期間を
より短いバックアップ間隔で実現できる．一方で，復旧時間や
スキャン停止による業務損失が大きい場合は，単に保持期間を延ばすだけでなく，
復旧手順の短縮，隔離・クローン・再デプロイの高速化，スキャン手順の
軽量化を合わせて検討すべきである．

この結論は，高頻度攻撃環境で特に使いやすい．固定費を含む最適間隔も，
攻撃頻度が高まるにつれて上の目安へ近づく．しかもその一次的な設計値は，
尾指数，攻撃頻度，検知レート，症状検知確率には依存しない．
これは，現場で不確かな攻撃分布を精密に推定し切れなくても，
全損時損失，復旧時間，停止損失，スキャン負荷，保持世代数という
運用上見積もりやすい量から，バックアップ間隔の初期設計を作れることを
意味する．もちろん，有限頻度での補正や許容区間の端点判定には
尾指数や検知率が影響するため，最終決定では第 5 節の目的関数を
端点も含めて評価する必要がある．

数学的な貢献としては，正の検知レート $\beta(I)>0$ のもとで
$\alpha\le1$ でも暦時間費用レートを定義できること，保持境界なしの
残余損失では $\kappa_s>0$ のもとで $\alpha=1$ を境界とする相転移が
現れること，早期検知失敗分岐の損失を指数傾き付き分岐確率
$\Hbeta$ と切断モーメント $\Mbeta$ から一貫して定義すると
上記の閉形式設計式が得られることを示した．

今後の課題は，早期検知分岐の費用を $\Delta$ 依存に拡張した場合の
補正項評価，有限 $\lambda$ での許容区間端点解の分類，復旧時間分布の
位相型分布への拡張，および実データからの尾指数推定である．
